\documentclass[a4paper,14pt]{extarticle}

% --- Кодировка и язык ---
\usepackage[utf8]{inputenc}
\usepackage[T2A]{fontenc}
\usepackage[english,russian]{babel}

% --- Шрифт: Tempora (Times-совместимый шрифт с поддержкой T2A и жирности) ---
\usepackage{tempora}

% --- Поля: левое 25 мм, правое 15 мм, верх/низ 20 мм ---
\RequirePackage[
  a4paper, includehead, includefoot, mag=1000,
  headsep=0mm, headheight=0mm,
  left=25mm, right=15mm, top=20mm, bottom=20mm
]{geometry}

% --- Полуторный межстрочный интервал ---
\usepackage{setspace}
\onehalfspacing

% --- Математика ---
\usepackage{amsmath,amssymb,amsthm}
\usepackage{mathtools}

% --- Таблицы, изображения ---
\usepackage{graphicx}
\usepackage{array}
\usepackage{longtable}
\usepackage{multirow}
\usepackage{enumitem}
\usepackage{indentfirst}
\usepackage{caption}
\usepackage{ragged2e}
\usepackage{pdfpages}

% --- Листинги кода ---
\usepackage{listings}
\usepackage{xcolor}

% Внутри lstlisting только ASCII: кириллица в listings + pdflatex
% даёт ошибку "TeX capacity exceeded [grouping levels=255]".
\lstdefinestyle{pythonstyle}{
  language=Python,
  basicstyle=\ttfamily\footnotesize,
  keywordstyle=\bfseries,
  commentstyle=\itshape\color{gray},
  stringstyle=\ttfamily,
  showstringspaces=false,
  breaklines=true,
  breakatwhitespace=false,
  columns=fullflexible,
  keepspaces=true,
  frame=single,
  framerule=0.4pt,
  numbers=left,
  numberstyle=\tiny\color{gray},
  numbersep=6pt,
  xleftmargin=1.5em,
  framexleftmargin=1.2em,
  tabsize=4,
  captionpos=b
}
\lstset{style=pythonstyle}

% --- Гиперссылки ---
\usepackage[hidelinks]{hyperref}

% --- Подписи ---
\DeclareCaptionLabelSeparator{emdash}{~--- }
\captionsetup[table]{
  labelsep=emdash,
  singlelinecheck=false,
  justification=raggedright,
  position=top,
  font=normalsize
}
\captionsetup[figure]{
  labelsep=emdash,
  singlelinecheck=false,
  justification=centering,
  position=bottom,
  font=normalsize
}
\captionsetup[lstlisting]{
  labelsep=emdash,
  singlelinecheck=false,
  justification=raggedright,
  font=normalsize
}

\usepackage{chngcntr}
\counterwithin{table}{section}
\counterwithin{figure}{section}
\counterwithin{equation}{section}

% --- Заголовки разделов ---
\usepackage{titlesec}
\titleformat{\section}{\normalfont\bfseries}{\thesection}{1em}{}
\titleformat{\subsection}{\normalfont\bfseries}{\thesubsection}{1em}{}
\titlespacing*{\section}{\parindent}{1em}{1em}
\titlespacing*{\subsection}{\parindent}{1em}{0.5em}

\setlength{\parindent}{1.25cm}

\newcommand{\struchead}[1]{%
  \clearpage
  \phantomsection
  \addcontentsline{toc}{section}{#1}%
  \begin{center}\textbf{\MakeUppercase{#1}}\end{center}\par\vspace{0.5em}
}

% --- Нумерация страниц снизу справа ---
\usepackage{fancyhdr}
\pagestyle{fancy}
\fancyhf{}
\fancyfoot[R]{\thepage}
\renewcommand{\headrulewidth}{0pt}
\renewcommand{\footrulewidth}{0pt}
\fancypagestyle{plain}{%
  \fancyhf{}%
  \fancyfoot[R]{\thepage}%
  \renewcommand{\headrulewidth}{0pt}%
}

% --- Содержание ---
\usepackage{tocloft}
\renewcommand{\cfttoctitlefont}{\hfill\bfseries\normalsize\MakeUppercase}
\renewcommand{\cftaftertoctitle}{\hfill}
\setcounter{tocdepth}{2}
\renewcommand{\contentsname}{Содержание}
\renewcommand{\cftsecaftersnum}{\hspace{0.5em}}
\renewcommand{\cftsubsecaftersnum}{\hspace{0.5em}}
\renewcommand{\cftsecfont}{\normalfont}
\renewcommand{\cftsubsecfont}{\normalfont}
\renewcommand{\cftsecpagefont}{\normalfont}
\renewcommand{\cftsubsecpagefont}{\normalfont}
\renewcommand{\cftsecdotsep}{\cftdotsep}
\renewcommand{\cftsubsecdotsep}{\cftdotsep}
\renewcommand{\cftsecleader}{\cftdotfill{\cftsecdotsep}}
\renewcommand{\cftsubsecleader}{\cftdotfill{\cftsubsecdotsep}}
% Отступы в содержании: задание (section) и подпункты (subsection)
\cftsetindents{section}{1.25cm}{2em}
\cftsetindents{subsection}{2.5cm}{3em}

\begin{document}

% ===================== ТИТУЛЬНЫЙ ЛИСТ =====================
% Если есть свой титульный PDF --- раскомментируйте строку ниже:
% \includepdf[pages={1}]{titulOtchet.pdf}

\begin{titlepage}
\centering
\vspace*{1cm}
{\large Саратовский национальный исследовательский государственный}\\[0.3em]
{\large университет имени Н.~Г.~Чернышевского}\\[0.5em]
{\large Механико-математический факультет}\\[0.3em]
{\large Кафедра математического и компьютерного моделирования}\\[2cm]

{\Large\bfseries Отчёт по учебной дисциплине}\\[0.5em]
{\Large\bfseries <<Основы объектно-ориентированного программирования>>}\\[1.5cm]

{\large Задание 1}\\[0.5em]
{\large Предметная область: космические миссии}\\[3cm]

\begin{flushright}
\begin{tabular}{l}
Студент: \underline{\hspace{5.5cm}}\\[1em]
Группа: \underline{\hspace{5.5cm}}\\[1em]
Преподаватель: Иванов~С.~В.
\end{tabular}
\end{flushright}

\vfill
Саратов --- \the\year
\end{titlepage}

% ===================== СОДЕРЖАНИЕ (отдельная страница) =====================
\setcounter{page}{2}
\tableofcontents
\thispagestyle{fancy}
\clearpage

% =====================================================================
% ЗАДАНИЕ 1  (далее будут \section{Задание 2. ...}, \section{Задание 3. ...})
% =====================================================================
\section{Задание 1. Классы предметной области}

\subsection{Постановка задачи}

В выбранной предметной области определить не менее трёх связанных между собой сущностей и их атрибуты. На языке Python описать классы, соответствующие найденным сущностям. Создать объекты описанной связки классов. Реализовать запросы, рекомендованные преподавателем. Для полученной программы построить диаграмму классов UML. При программировании каждый класс описывается в отдельном модуле (файле).

\subsection{Предметная область}

В качестве предметной области выбраны \textbf{космические миссии}. Необходимо хранить информацию о космонавтах, космических агентствах и миссиях.

В рассматриваемой предметной области выделены следующие классы:
\begin{itemize}[leftmargin=*,itemsep=2pt,topsep=4pt]
\item \textbf{космонавт} (\texttt{astronaut}) --- содержит информацию о космонавте и имеет следующие атрибуты:
  \begin{itemize}[leftmargin=*,itemsep=1pt]
  \item код (\texttt{code}) --- код, позволяющий однозначно определить космонавта, т.\,к.\ значения остальных атрибутов (ФИО) могут совпадать;
  \item фамилия (\texttt{surname});
  \item имя (\texttt{name});
  \item отчество (\texttt{secname});
  \item инициалы (\texttt{shortname}, \texttt{shortsecname}).
  \end{itemize}
\item \textbf{агентство} (\texttt{agency}) --- содержит информацию о космическом агентстве и имеет следующие атрибуты:
  \begin{itemize}[leftmargin=*,itemsep=1pt]
  \item код (\texttt{code});
  \item имя (\texttt{name}) --- полное название агентства;
  \item сокращённое обозначение (\texttt{shortname}), например <<Роскосмос>>, <<NASA>>.
  \end{itemize}
\item \textbf{миссия} (\texttt{mission}) --- содержит информацию о космической миссии и имеет следующие атрибуты:
  \begin{itemize}[leftmargin=*,itemsep=1pt]
  \item код (\texttt{code});
  \item имя (\texttt{name}) --- название миссии;
  \item длительность в днях (\texttt{days});
  \item год запуска (\texttt{year});
  \item эмблема / изображение (\texttt{img});
  \item агентство (\texttt{agency}) --- объект класса \texttt{agency};
  \item экипаж (\texttt{crew}) --- список объектов класса \texttt{astronaut}.
  \end{itemize}
\end{itemize}

Как видно из приведённого списка, все классы имеют сходные атрибуты --- \texttt{code} и \texttt{name}. В этом случае их можно выделить в отдельный класс-предок. Подобное выделение общих свойств называется \textbf{обобщением} (generalization). Назовём этот класс \texttt{general}.

\subsection{Класс general}

Ниже приведён листинг модуля \texttt{general.py}. Данный модуль включает описание класса \texttt{general}.

\begin{lstlisting}[caption={Модуль \texttt{general.py}}]
class general:
    def __init__(self, code=0, name=""):
        self.setCode(code)
        self.setName(name)

    def setCode(self, value):
        self.__code = value

    def setName(self, value):
        self.__name = value

    def getCode(self):
        return self.__code

    def getName(self):
        return self.__name
\end{lstlisting}

Описание класса:
\begin{lstlisting}
class general:
\end{lstlisting}

Далее описывается конструктор класса.
\begin{lstlisting}
def __init__(self, code=0, name=""):
    self.setCode(code)
    self.setName(name)
\end{lstlisting}

В данном случае конструктор позволяет при создании объекта задать начальные значения атрибутов. В заголовке функции перечислены входные параметры (\texttt{code}, \texttt{name}) и их значения по умолчанию. Первый параметр (\texttt{self}) является ссылкой объекта на самого себя, чтобы внутри методов обращаться к своим атрибутам и методам.

Поля и методы, название которых начинается с двойного подчёркивания, являются \textbf{приватными}. Если нужно сделать атрибуты доступными для просмотра и редактирования, описываются соответствующие открытые методы. Их принято называть по названию атрибута, добавляя слова \texttt{get} и \texttt{set}.

\begin{lstlisting}
def setCode(self, value):
    self.__code = value

def setName(self, value):
    self.__name = value
\end{lstlisting}

Get-методы возвращают значения атрибутов:
\begin{lstlisting}
def getCode(self):
    return self.__code

def getName(self):
    return self.__name
\end{lstlisting}

\subsection{Класс astronaut}

Далее опишем класс \texttt{astronaut} в модуле \texttt{astronaut.py}.

\begin{lstlisting}[caption={Модуль \texttt{astronaut.py}}]
from general import general

class astronaut(general):
    def __init__(self, code=0, surname="", name="", secname=""):
        general.__init__(self, code, name)
        self.setSurname(surname)
        self.setSecname(secname)

    def setSurname(self, value):
        self.__surname = value

    def setSecname(self, value):
        self.__secname = value

    def getSurname(self):
        return self.__surname

    def getSecname(self):
        return self.__secname

    def getShortname(self):
        if self.getName():
            return self.getName()[0]
        return ""

    def getShortsecname(self):
        if self.getSecname():
            return self.getSecname()[0]
        return ""

    def getInfoStr(self):
        s = self.getSurname()
        if self.getShortname():
            s += " %s." % (self.getShortname(),)
        if self.getShortsecname():
            s += " %s." % (self.getShortsecname(),)
        return s
\end{lstlisting}

В начале подключается модуль \texttt{general.py}, содержащий класс-предок:
\begin{lstlisting}
from general import general
\end{lstlisting}

Если класс наследуется от другого класса, название класса-предка указывается в скобках после имени наследуемого класса:
\begin{lstlisting}
class astronaut(general):
\end{lstlisting}

Чтобы класс унаследовал свойства предка, в конструкторе вызывается конструктор класса-предка:
\begin{lstlisting}
general.__init__(self, code, name)
\end{lstlisting}

Методы \texttt{getShortname()} и \texttt{getShortsecname()} возвращают инициалы имени и отчества. Метод \texttt{getInfoStr()} формирует краткую запись вида <<Гагарин Ю.~А.>>.

\subsection{Класс agency}

Аналогичным образом описываем класс \texttt{agency} в модуле \texttt{agency.py}.

\begin{lstlisting}[caption={Модуль \texttt{agency.py}}]
from general import general

class agency(general):
    def __init__(self, code=0, name="", shortname=""):
        general.__init__(self, code, name)
        self.setShortname(shortname)

    def setShortname(self, value):
        self.__shortname = value

    def getShortname(self):
        return self.__shortname
\end{lstlisting}

Класс \texttt{agency} наследует \texttt{code} и \texttt{name} от \texttt{general} и добавляет атрибут \texttt{shortname} --- краткое обозначение агентства.

\subsection{Класс generalList}

При описании класса \texttt{mission} обратим внимание на то, что миссия может иметь несколько космонавтов в экипаже, следовательно одним из атрибутов станет не один объект \texttt{astronaut}, а список таких объектов. Кроме того необходимы операции добавления и удаления космонавта и доступ по коду. Целесообразно выделить список в отдельный класс. В дальнейшем понадобятся аналогичные списки миссий и агентств, поэтому общие методы списков выделяются в класс-предок \texttt{generalList} (модуль \texttt{generalList.py}).

\begin{lstlisting}[caption={Модуль \texttt{generalList.py}}]
class generalList:
    def __init__(self):
        self.__list = []

    def clear(self):
        self.__list = []

    def findByCode(self, code):
        for item in self.__list:
            if item.getCode() == code:
                return item

    def getNewCode(self):
        codes = self.getCodes()
        if codes:
            return max(codes) + 1
        return 1

    def getCodes(self):
        return [item.getCode() for item in self.__list]

    def getItems(self):
        return list(self.__list)

    def appendItem(self, item):
        self.__list.append(item)

    def removeItem(self, code):
        item = self.findByCode(code)
        if item is not None:
            self.__list.remove(item)
            return True
        return False
\end{lstlisting}

Объекты хранятся в приватном списке \texttt{\_\_list}, который инициализируется в конструкторе. Метод \texttt{clear()} очищает список. Метод \texttt{findByCode(code)} возвращает ссылку на объект по значению атрибута \texttt{code}. Метод \texttt{getNewCode()} возвращает следующий свободный код. Метод \texttt{getCodes()} возвращает список кодов объектов. Методы \texttt{appendItem} и \texttt{removeItem} добавляют и удаляют элементы.

\subsection{Класс astronautList}

От класса \texttt{generalList} наследуем класс \texttt{astronautList}, который содержит список космонавтов и функции доступа к ним (модуль \texttt{astronautList.py}).

\begin{lstlisting}[caption={Модуль \texttt{astronautList.py}}]
from astronaut import astronaut
from generalList import generalList

class astronautList(generalList):
    def appendItem(self, item):
        if not isinstance(item, astronaut):
            raise TypeError("astronautList accepts only astronaut")
        generalList.appendItem(self, item)

    def createItem(self, code, surname="", name="", secname=""):
        item = astronaut(code, surname, name, secname)
        self.appendItem(item)
        return item

    def newItem(self, surname="", name="", secname=""):
        return self.createItem(self.getNewCode(), surname, name, secname)

    def getInfoStr(self):
        return ", ".join(item.getInfoStr() for item in self.getItems())
\end{lstlisting}

Метод \texttt{appendItem} переопределён: в список можно добавлять только объекты класса \texttt{astronaut}. Метод \texttt{createItem} создаёт космонавта с заданным кодом, \texttt{newItem} --- с автоматически сгенерированным кодом. Метод \texttt{getInfoStr()} возвращает строку экипажа через запятую, например <<Гагарин Ю.~А., Титов Г.~С.>>.

\subsection{Класс mission}

Далее описываем класс \texttt{mission} в модуле \texttt{mission.py}.

\begin{lstlisting}[caption={Модуль \texttt{mission.py}}]
from astronautList import astronautList
from general import general

class mission(general):
    def __init__(self, code=0, name="", agency=None,
                 year=0, days=0, img=""):
        general.__init__(self, code, name)
        self.__crew = astronautList()
        self.setAgency(agency)
        self.setDays(days)
        self.setYear(year)
        self.setImg(img)

    def setImg(self, value):
        self.__img = value

    def setAgency(self, value):
        self.__agency = value

    def setDays(self, value):
        self.__days = value

    def setYear(self, value):
        self.__year = value

    def getImg(self):
        return self.__img

    def getAgency(self):
        return self.__agency

    def getAgencyCode(self):
        if self.__agency:
            return self.__agency.getCode()

    def getAgencyName(self):
        if self.__agency:
            return self.__agency.getName()
        return ""

    def getAgencyShortName(self):
        if self.__agency:
            return self.__agency.getShortname()
        return ""

    def getDays(self):
        return self.__days

    def getYear(self):
        return self.__year

    def appendAstronaut(self, item):
        self.__crew.appendItem(item)

    def removeAstronaut(self, code):
        return self.__crew.removeItem(code)

    def clearCrew(self):
        self.__crew.clear()

    def getAstronaut(self, code):
        return self.__crew.findByCode(code)

    def getCrew(self):
        return self.__crew

    def getAstronautCodes(self):
        return self.__crew.getCodes()

    def getCrewInfoStr(self):
        return self.__crew.getInfoStr()

    def getInfoStr(self):
        s = self.getCrewInfoStr()
        if s:
            s += " "
        s += "%s - %s, %s. - %s d." % (
            self.getName(),
            self.getAgencyShortName(),
            str(self.getYear()),
            str(self.getDays()),
        )
        return s
\end{lstlisting}

В конструкторе используется конструктор класса-предка. В качестве атрибута экипажа используется объект класса \texttt{astronautList}. Значение по умолчанию для ссылки на агентство задаётся как \texttt{None}.

Метод \texttt{getInfoStr()} возвращает сводку по миссии в формате:
\begin{center}
\textit{Гагарин Ю.~А. Восток-1 --- Роскосмос, 1961. --- 1 дн.}
\end{center}

\subsection{Демонстрационная программа}

Программа в файле \texttt{main.py} создаёт объекты описанных классов и демонстрирует работу запросов.

\begin{lstlisting}[caption={Фрагмент модуля \texttt{main.py}: создание объектов}]
from agency import agency
from astronaut import astronaut
from mission import mission

roscosmos = agency(1, "Roscosmos", "Roscosmos")
nasa = agency(2, "National Aeronautics and Space Administration", "NASA")

gagarin = astronaut(1, "Gagarin", "Yuri", "Alekseevich")
armstrong = astronaut(3, "Armstrong", "Neil", "Alden")
aldrin = astronaut(4, "Aldrin", "Buzz", "")

vostok1 = mission(1, "Vostok-1", roscosmos, year=1961, days=1,
                  img="vostok1.png")
vostok1.appendAstronaut(gagarin)

apollo11 = mission(2, "Apollo 11", nasa, year=1969, days=8,
                   img="apollo11.png")
apollo11.appendAstronaut(armstrong)
apollo11.appendAstronaut(aldrin)

print(vostok1.getInfoStr())
print(apollo11.getInfoStr())
\end{lstlisting}

Создаются объекты классов \texttt{agency} и \texttt{astronaut}, затем объекты \texttt{mission}; космонавты добавляются в экипаж методом \texttt{appendAstronaut}. Метод \texttt{getInfoStr()} выводит сводку по миссии.

Результат выполнения программы (фрагмент):
\begin{quote}
\ttfamily\footnotesize
Gagarin Y. A. Vostok-1 - Roscosmos, 1961. - 1 d.\\
Armstrong N. A., Aldrin B. Apollo 11 - NASA, 1969. - 8 d.
\end{quote}

\noindent В реальном коде программы используются русские строки (<<Гагарин Ю.~А. Восток-1 --- Роскосмос, 1961. --- 1 дн.>> и~т.\,д.); в листингах отчёта кириллица заменена на латиницу из\-за ограничений пакета \texttt{listings} при компиляции pdfLaTeX.

Реализованы также запросы: поиск миссии по коду, вывод миссий выбранного агентства, поиск миссий с участием заданного космонавта, добавление и удаление членов экипажа, генерация нового кода через \texttt{getNewCode()} / \texttt{newItem()}.

\subsection{Диаграмма классов UML}

Диаграмма классов полученной программы приведена на рисунке~\ref{fig:uml1}.

% ============================================================
% ВСТАВЬТЕ СЮДА КАРТИНКУ UML:
% 1) Сгенерируйте диаграмму на plantuml.com (скрипт diagram.puml)
% 2) Сохраните файл как report/diagram1.png (или .pdf)
% 3) При необходимости поправьте имя файла ниже
% ============================================================
\begin{figure}[ht]
\centering
\includegraphics[width=0.95\textwidth]{diagram1.png}
\caption{Диаграмма классов для задания 1}
\label{fig:uml1}
\end{figure}

Скрипт PlantUML для диаграммы классов содержится в файле \texttt{diagram.puml} проекта.

\subsection{Выводы}

В ходе выполнения задания~1 для предметной области <<космические миссии>> выделены связанные сущности \texttt{astronaut}, \texttt{agency} и \texttt{mission}. Общие атрибуты вынесены в класс-предок \texttt{general}, общие операции над коллекциями --- в \texttt{generalList}. Реализованы наследование, инкапсуляция и агрегация (миссия содержит экипаж и ссылку на агентство). Созданы объекты связки классов, выполнены типовые запросы, построена диаграмма классов UML.

% =====================================================================
% ЗАДАНИЕ 2
% =====================================================================
\clearpage
\section{Задание 2. Чтение и запись данных в формате XML}

\subsection{Постановка задачи}

Написать функции для считывания и записи данных о выбранных сущностях из файла в формате XML, используя при этом классы, полученные при выполнении задания~1. Для полученной программы построить диаграмму классов UML. Реализовать запросы, рекомендованные преподавателем. Для демонстрации работы функции организовать вывод данных на экран в консоли.

\subsection{Исходный XML-файл}

В текстовом редакторе создан файл \texttt{old.xml}. В нём записаны данные о космонавтах, агентствах и миссиях в формате XML (ниже --- латинская версия листинга; в рабочем файле проекта используются русские строки).

\begin{lstlisting}[language={},caption={Файл \texttt{old.xml} (фрагмент структуры)}]
<?xml version="1.0" encoding="utf-8"?>
<cosmos>
  <astronaut code="1" surname="Gagarin" name="Yuri"
             secname="Alekseevich"/>
  <astronaut code="3" surname="Armstrong" name="Neil"
             secname="Alden"/>
  <astronaut code="4" surname="Aldrin" name="Buzz" secname=""/>
  <agency code="1" name="Roscosmos" shortname="Roscosmos"/>
  <agency code="2" name="NASA" shortname="NASA"/>
  <mission code="1" name="Vostok-1" img="vostok1.png"
           year="1961" days="1" agency="1">
    <astronaut code="1"/>
  </mission>
  <mission code="2" name="Apollo 11" img="apollo11.png"
           year="1969" days="8" agency="2">
    <astronaut code="3"/>
    <astronaut code="4"/>
  </mission>
</cosmos>
\end{lstlisting}

Корневой элемент --- \texttt{cosmos}. Внутри него перечисляются элементы \texttt{astronaut}, \texttt{agency} и \texttt{mission}. У миссии вложенные элементы \texttt{astronaut} задают экипаж по кодам; атрибут \texttt{agency} ссылается на код агентства.

\subsection{Списки agencyList и missionList}

Для хранения считанных данных потребуются списки космонавтов, агентств и миссий. Список космонавтов уже реализован в задании~1 (\texttt{astronautList}). Добавим модули \texttt{agencyList.py} и \texttt{missionList.py}.

\begin{lstlisting}[caption={Модуль \texttt{agencyList.py}}]
from agency import agency
from generalList import generalList

class agencyList(generalList):
    def appendItem(self, value):
        if isinstance(value, agency):
            generalList.appendItem(self, value)

    def newItem(self, name, shortname=""):
        a = agency(self.getNewCode(), name, shortname)
        self.appendItem(a)
        return a

    def createItem(self, code, name, shortname=""):
        if code in self.getCodes():
            print("Agency with code %s already exists" % (code,))
        else:
            a = agency(code, name, shortname)
            self.appendItem(a)
            return a
\end{lstlisting}

\begin{lstlisting}[caption={Модуль \texttt{missionList.py}}]
from generalList import generalList
from mission import mission

class missionList(generalList):
    def appendItem(self, value):
        if isinstance(value, mission):
            generalList.appendItem(self, value)

    def createItem(self, code, name, agency=None,
                   year=0, days=0, img=""):
        if code in self.getCodes():
            print("Mission with code %s already exists" % (code,))
        else:
            m = mission(code, name, agency, year, days, img)
            self.appendItem(m)
            return m

    def newItem(self, name, agency=None, year=0, days=0, img=""):
        m = mission(self.getNewCode(), name, agency, year, days, img)
        self.appendItem(m)
        return m
\end{lstlisting}

Метод \texttt{appendItem} принимает только объекты нужного типа. Методы \texttt{createItem} и \texttt{newItem} создают элемент с заданным или автоматически сгенерированным кодом.

\subsection{Класс cosmos}

Все списки объектов хранятся в объекте класса \texttt{cosmos} (аналог класса \texttt{library} из методички). Класс также содержит функции создания и удаления объектов.

\begin{lstlisting}[caption={Модуль \texttt{cosmos.py} (основные методы)}]
from agencyList import agencyList
from astronautList import astronautList
from missionList import missionList

class cosmos:
    def __init__(self):
        self.__astronautList = astronautList()
        self.__agencyList = agencyList()
        self.__missionList = missionList()

    def clear(self):
        self.__missionList.clear()
        self.__astronautList.clear()
        self.__agencyList.clear()

    def createAstronaut(self, code, surname, name="", secname=""):
        return self.__astronautList.createItem(
            code, surname, name, secname)

    def getAstronaut(self, code):
        return self.__astronautList.findByCode(code)

    def getAstronautList(self):
        return self.__astronautList.getItems()

    def getAstronautCodes(self):
        return self.__astronautList.getCodes()

    def createAgency(self, code, name, shortname=""):
        return self.__agencyList.createItem(code, name, shortname)

    def getAgency(self, code):
        return self.__agencyList.findByCode(code)

    def getAgencyList(self):
        return self.__agencyList.getItems()

    def createMission(self, code, name, agency=None,
                      year=0, days=0, img=""):
        return self.__missionList.createItem(
            code, name, agency, year, days, img)

    def getMission(self, code):
        return self.__missionList.findByCode(code)

    def getMissionList(self):
        return self.__missionList.getItems()

    def getMissionCodes(self):
        return self.__missionList.getCodes()
\end{lstlisting}

При удалении космонавта он также убирается из экипажей всех миссий; при удалении агентства у связанных миссий ссылка на агентство обнуляется (\texttt{setAgency(None)}).

\subsection{Класс data}

С учётом последующих заданий (JSON, SQL) реализован класс \texttt{data} --- общий предок для всех классов чтения и записи. Различаться наследники будут только методами \texttt{read()} и \texttt{write()}, зависящими от формата.

\begin{lstlisting}[caption={Модуль \texttt{data.py}}]
class data:
    def __init__(self, lib=None, inp="", out=""):
        self.setLib(lib)
        self.setInp(inp)
        self.setOut(out)

    def setLib(self, value):
        self.__lib = value

    def setInp(self, value):
        self.__inp = value

    def setOut(self, value):
        self.__out = value

    def getLib(self):
        return self.__lib

    def getInp(self):
        return self.__inp

    def getOut(self):
        return self.__out

    def readFile(self, lib=None, filename=""):
        if lib:
            self.setLib(lib)
        if filename:
            self.setInp(filename)
        self.read()

    def writeFile(self, lib=None, filename=""):
        if lib:
            self.setLib(lib)
        if filename:
            self.setOut(filename)
        self.write()

    def read(self):
        pass

    def write(self):
        pass
\end{lstlisting}

Атрибуты: \texttt{lib} --- ссылка на каталог \texttt{cosmos}, \texttt{inp} / \texttt{out} --- имена входного и выходного файлов. Методы \texttt{read()} и \texttt{write()} в базовом классе пустые (заглушки).

\subsection{Класс dataxml}

Класс \texttt{dataxml} наследуется от \texttt{data}; в нём переопределены методы \texttt{read()} и \texttt{write()} для работы с XML через модуль \texttt{xml.dom.minidom}.

\begin{lstlisting}[caption={Модуль \texttt{dataxml.py}: метод \texttt{read()}}]
import xml.dom.minidom
from data import data

class dataxml(data):
    def read(self):
        dom = xml.dom.minidom.parse(self.getInp())
        dom.normalize()
        for node in dom.childNodes[0].childNodes:
            if (node.nodeType == node.ELEMENT_NODE) and \
               (node.nodeName == "astronaut"):
                code, surname, name, secname = 0, "", "", ""
                for t in node.attributes.items():
                    if t[0] == "code":
                        code = int(t[1])
                    if t[0] == "name":
                        name = t[1]
                    if t[0] == "surname":
                        surname = t[1]
                    if t[0] == "secname":
                        secname = t[1]
                self.getLib().createAstronaut(
                    code, surname, name, secname)

            if (node.nodeType == node.ELEMENT_NODE) and \
               (node.nodeName == "agency"):
                code, name, shortname = 0, "", ""
                for t in node.attributes.items():
                    if t[0] == "code":
                        code = int(t[1])
                    if t[0] == "name":
                        name = t[1]
                    if t[0] == "shortname":
                        shortname = t[1]
                self.getLib().createAgency(code, name, shortname)

            if (node.nodeType == node.ELEMENT_NODE) and \
               (node.nodeName == "mission"):
                code, name, img = 0, "", ""
                agency, year, days = None, 0, 0
                for t in node.attributes.items():
                    if t[0] == "code":
                        code = int(t[1])
                    if t[0] == "name":
                        name = t[1]
                    if t[0] == "img":
                        img = t[1]
                    if t[0] == "year":
                        year = int(t[1])
                    if t[0] == "days":
                        days = int(t[1])
                    if t[0] == "agency":
                        agency = self.getLib().getAgency(int(t[1]))
                m = self.getLib().createMission(
                    code, name, agency, year, days, img)
                for n in node.childNodes:
                    if (n.nodeType == n.ELEMENT_NODE) and \
                       (n.nodeName == "astronaut"):
                        for t in n.attributes.items():
                            if t[0] == "code":
                                person = self.getLib().getAstronaut(
                                    int(t[1]))
                                if person and m:
                                    m.appendAstronaut(person)
\end{lstlisting}

Сначала читаются космонавты и агентства, затем миссии: для миссии по коду находится объект агентства, создаётся миссия, в экипаж добавляются космонавты по вложенным элементам \texttt{astronaut}.

\begin{lstlisting}[caption={Модуль \texttt{dataxml.py}: метод \texttt{write()}}]
    def write(self):
        dom = xml.dom.minidom.Document()
        root = dom.createElement("cosmos")
        dom.appendChild(root)

        for a in self.getLib().getAstronautList():
            aut = dom.createElement("astronaut")
            aut.setAttribute("code", str(a.getCode()))
            aut.setAttribute("surname", a.getSurname())
            aut.setAttribute("name", a.getName())
            aut.setAttribute("secname", a.getSecname())
            root.appendChild(aut)

        for p in self.getLib().getAgencyList():
            ag = dom.createElement("agency")
            ag.setAttribute("code", str(p.getCode()))
            ag.setAttribute("name", p.getName())
            ag.setAttribute("shortname", p.getShortname())
            root.appendChild(ag)

        for m in self.getLib().getMissionList():
            ms = dom.createElement("mission")
            ms.setAttribute("code", str(m.getCode()))
            ms.setAttribute("name", m.getName())
            ms.setAttribute("img", m.getImg())
            ms.setAttribute("year", str(m.getYear()))
            ms.setAttribute("days", str(m.getDays()))
            agency_code = m.getAgencyCode()
            ms.setAttribute(
                "agency",
                str(agency_code) if agency_code is not None else "")
            for ac in m.getAstronautCodes():
                aut = dom.createElement("astronaut")
                aut.setAttribute("code", str(ac))
                ms.appendChild(aut)
            root.appendChild(ms)

        with open(self.getOut(), "w", encoding="utf-8") as f:
            f.write(dom.toprettyxml())
\end{lstlisting}

Метод \texttt{write()} строит XML-документ, заполняет его данными из каталога и сохраняет в файл \texttt{out} с кодировкой UTF-8.

\subsection{Демонстрационная программа}

Программа в файле \texttt{main2.py} создаёт объекты классов \texttt{cosmos} и \texttt{dataxml}, читает \texttt{old.xml}, выводит данные в консоль и записывает \texttt{new.xml}.

\begin{lstlisting}[caption={Модуль \texttt{main2.py}}]
from cosmos import cosmos
from dataxml import dataxml

catalog = cosmos()
dat = dataxml(catalog, "old.xml", "new.xml")
dat.read()

for m in catalog.getMissionList():
    print(m.getInfoStr())

dat.write()
\end{lstlisting}

Импортируются необходимые классы, создаётся каталог и объект \texttt{dataxml} с указанием входного и выходного файлов. Вызов \texttt{read()} заполняет каталог; цикл по миссиям выводит сводки; \texttt{write()} сохраняет данные в \texttt{new.xml}.

Результат выполнения (фрагмент):
\begin{quote}
\ttfamily\footnotesize
Gagarin Y. A. Vostok-1 - Roscosmos, 1961. - 1 d.\\
Armstrong N. A., Aldrin B. Apollo 11 - NASA, 1969. - 8 d.\\
Tereshkova V. V. Vostok-6 - Roscosmos, 1963. - 3 d.
\end{quote}

Дополнительно реализованы запросы: вывод миссий выбранного агентства, поиск космонавта по коду, контрольное перечитывание записанного файла \texttt{new.xml}.

\subsection{Диаграмма классов UML}

Диаграмма классов для заданий~2 (и последующих форматов данных) приведена на рисунке~\ref{fig:uml2}. По сравнению с заданием~1 добавлены классы \texttt{agencyList}, \texttt{missionList}, \texttt{cosmos}, \texttt{data}, \texttt{dataxml}.

% ============================================================
% ВСТАВЬТЕ СЮДА КАРТИНКУ UML ЗАДАНИЯ 2:
% 1) Сгенерируйте диаграмму на plantuml.com (скрипт diagram2.puml)
% 2) Сохраните файл как report/diagram2.png
% ============================================================
\begin{figure}[ht]
\centering
\includegraphics[width=0.95\textwidth]{diagram2.png}
\caption{Диаграмма классов для задания 2}
\label{fig:uml2}
\end{figure}

Скрипт PlantUML содержится в файле \texttt{diagram2.puml} проекта.

\subsection{Выводы}

В задании~2 на основе классов задания~1 реализованы контейнер \texttt{cosmos}, специализированные списки и иерархия \texttt{data}~/~\texttt{dataxml} для чтения и записи XML. Данные из \texttt{old.xml} загружаются в объекты, обрабатываются запросами и сохраняются в \texttt{new.xml}. Построена диаграмма классов UML.

% =====================================================================
% ЗАДАНИЕ 2a
% =====================================================================
\clearpage
\section{Задание 2a. Чтение и запись данных в формате JSON}

\subsection{Постановка задачи}

Написать функции для считывания и записи данных о выбранных сущностях из файла в формате JSON, используя при этом классы, полученные при выполнении задания~1. Для полученной программы построить диаграмму классов UML. Реализовать запросы, рекомендованные преподавателем. Для демонстрации работы функции организовать вывод данных на экран в консоли.

\subsection{Исходный JSON-файл}

В текстовом редакторе создан файл \texttt{old.json}. Структура данных аналогична XML-каталогу: массивы космонавтов, агентств и миссий (ниже --- латинская версия листинга; в рабочем файле проекта используются русские строки).

\begin{lstlisting}[language={},caption={Файл \texttt{old.json} (фрагмент структуры)}]
{
  "astronauts": [
    {"code": 1, "surname": "Gagarin", "name": "Yuri",
     "secname": "Alekseevich"},
    {"code": 3, "surname": "Armstrong", "name": "Neil",
     "secname": "Alden"},
    {"code": 4, "surname": "Aldrin", "name": "Buzz",
     "secname": ""}
  ],
  "agencies": [
    {"code": 1, "name": "Roscosmos", "shortname": "Roscosmos"},
    {"code": 2, "name": "NASA", "shortname": "NASA"}
  ],
  "missions": [
    {
      "code": 1,
      "name": "Vostok-1",
      "img": "vostok1.png",
      "year": 1961,
      "days": 1,
      "agency": 1,
      "astronauts": [1]
    },
    {
      "code": 2,
      "name": "Apollo 11",
      "img": "apollo11.png",
      "year": 1969,
      "days": 8,
      "agency": 2,
      "astronauts": [3, 4]
    }
  ]
}
\end{lstlisting}

Ключи верхнего уровня: \texttt{astronauts}, \texttt{agencies}, \texttt{missions}. У миссии поле \texttt{agency} --- код агентства, поле \texttt{astronauts} --- список кодов членов экипажа.

\subsection{Классы cosmos и data}

Класс \texttt{cosmos} и базовый класс \texttt{data} используются те же, что в задании~2. Дополнительно описывается только наследник \texttt{datajson} для формата JSON.

\subsection{Класс datajson}

Класс \texttt{datajson} наследуется от \texttt{data}; в нём переопределены методы \texttt{read()} и \texttt{write()} с использованием модуля \texttt{json}.

\begin{lstlisting}[caption={Модуль \texttt{datajson.py}: метод \texttt{read()}}]
import json
from data import data

class datajson(data):
    def read(self):
        with open(self.getInp(), "r", encoding="utf-8") as f:
            data = json.load(f)
        for k in data.keys():
            if k == "astronauts":
                for a in data[k]:
                    code, surname, name, secname = 0, "", "", ""
                    for ak in a.keys():
                        if ak == "code":
                            code = a[ak]
                        if ak == "name":
                            name = a[ak]
                        if ak == "surname":
                            surname = a[ak]
                        if ak == "secname":
                            secname = a[ak]
                    self.getLib().createAstronaut(
                        code, surname, name, secname)
            if k == "agencies":
                for a in data[k]:
                    code, name, shortname = 0, "", ""
                    for ak in a.keys():
                        if ak == "code":
                            code = a[ak]
                        if ak == "name":
                            name = a[ak]
                        if ak == "shortname":
                            shortname = a[ak]
                    self.getLib().createAgency(code, name, shortname)
            if k == "missions":
                for a in data[k]:
                    code, name, img = 0, "", ""
                    agency, year, days, crew = None, 0, 0, []
                    for ak in a.keys():
                        if ak == "code":
                            code = a[ak]
                        if ak == "name":
                            name = a[ak]
                        if ak == "img":
                            img = a[ak]
                        if ak == "year":
                            year = a[ak]
                        if ak == "days":
                            days = a[ak]
                        if ak == "agency":
                            agency = self.getLib().getAgency(int(a[ak]))
                        if ak == "astronauts":
                            crew = a[ak]
                    m = self.getLib().createMission(
                        code, name, agency, year, days, img)
                    for ac in crew:
                        person = self.getLib().getAstronaut(ac)
                        if person and m:
                            m.appendAstronaut(person)
\end{lstlisting}

Файл читается через \texttt{json.load}. Для каждого ключа создаются соответствующие объекты каталога. Для миссий сначала создаётся объект, затем в экипаж добавляются космонавты по списку кодов.

\begin{lstlisting}[caption={Модуль \texttt{datajson.py}: метод \texttt{write()}}]
    def write(self):
        data = {"astronauts": [], "agencies": [], "missions": []}
        for a in self.getLib().getAstronautList():
            ad = {}
            ad["code"] = a.getCode()
            ad["surname"] = a.getSurname()
            ad["name"] = a.getName()
            ad["secname"] = a.getSecname()
            data["astronauts"].append(ad)
        for p in self.getLib().getAgencyList():
            pd = {}
            pd["code"] = p.getCode()
            pd["name"] = p.getName()
            pd["shortname"] = p.getShortname()
            data["agencies"].append(pd)
        for m in self.getLib().getMissionList():
            md = {}
            md["code"] = m.getCode()
            md["name"] = m.getName()
            md["img"] = m.getImg()
            md["year"] = m.getYear()
            md["days"] = m.getDays()
            md["agency"] = m.getAgencyCode()
            md["astronauts"] = m.getAstronautCodes()
            data["missions"].append(md)
        with open(self.getOut(), "w", encoding="utf-8") as f:
            json.dump(data, f, indent=2, ensure_ascii=False)
\end{lstlisting}

Метод \texttt{write()} формирует словарь той же структуры и сохраняет его через \texttt{json.dump} с параметрами \texttt{indent=2} и \texttt{ensure\_ascii=False} (сохранение кириллицы без экранирования).

\subsection{Демонстрационная программа}

Программа в файле \texttt{main2a.py} создаёт объекты классов \texttt{cosmos} и \texttt{datajson}, читает \texttt{old.json}, выводит данные в консоль и записывает \texttt{new.json}.

\begin{lstlisting}[caption={Модуль \texttt{main2a.py}}]
from cosmos import cosmos
from datajson import datajson

catalog = cosmos()
dat = datajson(catalog, "old.json", "new.json")
dat.read()

for m in catalog.getMissionList():
    print(m.getInfoStr())

dat.write()
\end{lstlisting}

Результат выполнения (фрагмент):
\begin{quote}
\ttfamily\footnotesize
Gagarin Y. A. Vostok-1 - Roscosmos, 1961. - 1 d.\\
Armstrong N. A., Aldrin B. Apollo 11 - NASA, 1969. - 8 d.\\
Tereshkova V. V. Vostok-6 - Roscosmos, 1963. - 3 d.
\end{quote}

Дополнительно реализованы запросы: вывод миссий выбранного агентства, поиск космонавта по коду, контрольное перечитывание \texttt{new.json}.

\subsection{Диаграмма классов UML}

Диаграмма классов для задания~2a приведена на рисунке~\ref{fig:uml2a}. По сравнению с заданием~2 добавлен класс \texttt{datajson} (рядом с \texttt{dataxml} как наследник \texttt{data}).

% ============================================================
% ВСТАВЬТЕ СЮДА КАРТИНКУ UML ЗАДАНИЯ 2a:
% 1) Сгенерируйте диаграмму на plantuml.com (скрипт diagram2a.puml)
% 2) Сохраните файл как report/diagram2a.png
% ============================================================
\begin{figure}[ht]
\centering
\includegraphics[width=0.95\textwidth]{diagram2a.png}
\caption{Диаграмма классов для задания 2a}
\label{fig:uml2a}
\end{figure}

Скрипт PlantUML содержится в файле \texttt{diagram2a.puml} проекта.

\subsection{Выводы}

В задании~2a реализован класс \texttt{datajson} для чтения и записи каталога космических миссий в формате JSON с использованием классов заданий~1 и~2. Данные из \texttt{old.json} загружаются в объекты, обрабатываются запросами и сохраняются в \texttt{new.json}. Построена диаграмма классов UML.

% =====================================================================
% ЗАДАНИЕ 3
% =====================================================================
\clearpage
\section{Задание 3. Чтение и запись данных в SQLite}

\subsection{Постановка задачи}

Написать функции для считывания и записи данных о выбранных сущностях из базы данных SQLite, используя при этом классы, полученные при выполнении заданий~1 и~2. Реализовать запросы, рекомендованные преподавателем. Для демонстрации работы функции организовать вывод данных на экран в консоли.

\subsection{Схема базы данных}

Для хранения каталога космических миссий используются четыре таблицы. Связь <<многие ко многим>> между миссиями и космонавтами реализуется через таблицу \texttt{mission\_astronaut}.

\begin{lstlisting}[language=SQL,caption={Схема SQLite}]
PRAGMA foreign_keys = ON;

CREATE TABLE astronaut (
  code integer primary key,
  surname text,
  name text,
  secname text
);

CREATE TABLE agency (
  code integer primary key,
  name text,
  shortname text
);

CREATE TABLE mission (
  code integer primary key,
  name text,
  img text,
  year integer,
  days integer,
  agency integer references agency(code)
    on update cascade on delete set null
);

CREATE TABLE mission_astronaut (
  code integer primary key autoincrement,
  mission integer references mission(code)
    on update cascade on delete cascade,
  astronaut integer references astronaut(code)
    on update cascade on delete cascade,
  unique(mission, astronaut)
);
\end{lstlisting}

ER-диаграмма приведена на рисунке~\ref{fig:er3}.

% ============================================================
% ВСТАВЬТЕ СЮДА ER-ДИАГРАММУ:
% 1) Сгенерируйте на plantuml.com (скрипт diagram3_er.puml)
% 2) Сохраните как report/diagram3_er.png
% ============================================================
\begin{figure}[ht]
\centering
\includegraphics[width=0.85\textwidth]{diagram3_er.png}
\caption{ER-диаграмма базы данных SQLite}
\label{fig:er3}
\end{figure}

\subsection{Класс datasql}

Класс \texttt{datasql} наследуется от \texttt{data}; в нём переопределены методы \texttt{read()} и \texttt{write()} с использованием модуля \texttt{sqlite3}.

\begin{lstlisting}[caption={Модуль \texttt{datasql.py}: метод \texttt{read()}}]
import sqlite3 as db
from data import data

class datasql(data):
    def read(self):
        conn = db.connect(self.getInp())
        curs = conn.cursor()
        curs.execute(
            "select code,surname,name,secname from astronaut")
        for r in curs.fetchall():
            self.getLib().createAstronaut(r[0], r[1], r[2], r[3])

        curs.execute("select code,name,shortname from agency")
        for r in curs.fetchall():
            self.getLib().createAgency(r[0], r[1], r[2])

        curs.execute(
            "select code,name,img,year,days,agency from mission")
        for r in curs.fetchall():
            agency = None
            if r[5] is not None:
                agency = self.getLib().getAgency(int(r[5]))
            self.getLib().createMission(
                r[0], r[1], agency, r[3], r[4], r[2])

        curs.execute(
            "select mission,astronaut from mission_astronaut")
        for r in curs.fetchall():
            m = self.getLib().getMission(r[0])
            person = self.getLib().getAstronaut(r[1])
            if m and person:
                m.appendAstronaut(person)
        conn.close()
\end{lstlisting}

Сначала загружаются космонавты и агентства, затем миссии (со ссылкой на агентство), затем связи экипажа из \texttt{mission\_astronaut}.

\begin{lstlisting}[caption={Модуль \texttt{datasql.py}: метод \texttt{write()}}]
    def write(self):
        conn = db.connect(self.getOut())
        curs = conn.cursor()
        curs.executescript(emptydb)

        for a in self.getLib().getAstronautList():
            curs.execute(
                "insert into astronaut(code,surname,name,secname)"
                " values(?,?,?,?)",
                (a.getCode(), a.getSurname(),
                 a.getName(), a.getSecname()))

        for p in self.getLib().getAgencyList():
            curs.execute(
                "insert into agency(code,name,shortname)"
                " values(?,?,?)",
                (p.getCode(), p.getName(), p.getShortname()))

        for m in self.getLib().getMissionList():
            curs.execute(
                "insert into mission(code,name,img,year,days,agency)"
                " values(?,?,?,?,?,?)",
                (m.getCode(), m.getName(), m.getImg(),
                 m.getYear(), m.getDays(), m.getAgencyCode()))
            for ac in m.getAstronautCodes():
                curs.execute(
                    "insert into mission_astronaut"
                    "(mission,astronaut) values(?,?)",
                    (m.getCode(), ac))

        conn.commit()
        conn.close()
\end{lstlisting}

Перед вставкой данных выполняется скрипт \texttt{emptydb}, создающий пустую схему таблиц. Для проверки достаточно иметь исходный XML/JSON: база создаётся программно при записи.

\subsection{Демонстрационная программа}

Программа \texttt{main3.py} проверяет круг обмена данными: XML $\to$ SQLite $\to$ объекты $\to$ XML.

\begin{lstlisting}[caption={Модуль \texttt{main3.py}}]
import os
from cosmos import cosmos
from dataxml import dataxml
from datasql import datasql

cat1 = cosmos()
cat2 = cosmos()
dxml1 = dataxml(cat1, "old.xml", "new.xml")
dxml2 = dataxml(cat2, "old.xml", "roundtrip.xml")
dsql1 = datasql(cat1, "new.sqlite", "new.sqlite")
dsql2 = datasql(cat2, "new.sqlite", "new.sqlite")

dxml1.read()
if os.path.isfile(dsql1.getOut()):
    os.remove(dsql1.getOut())
dsql1.write()
dsql2.read()

for m in cat2.getMissionList():
    print(m.getInfoStr())

dxml2.write()
\end{lstlisting}

Сначала данные читаются из \texttt{old.xml} в первый каталог, записываются в \texttt{new.sqlite}, затем читаются во второй каталог и выводятся в консоль; дополнительно сохраняется \texttt{roundtrip.xml}.

Результат выполнения (фрагмент):
\begin{quote}
\ttfamily\footnotesize
Gagarin Y. A. Vostok-1 - Roscosmos, 1961. - 1 d.\\
Armstrong N. A., Aldrin B. Apollo 11 - NASA, 1969. - 8 d.\\
Tereshkova V. V. Vostok-6 - Roscosmos, 1963. - 3 d.
\end{quote}

Дополнительно реализованы запросы: вывод миссий выбранного агентства, контроль числа сущностей после круга обмена.

\subsection{Диаграмма классов UML}

Диаграмма классов для задания~3 приведена на рисунке~\ref{fig:uml3}. По сравнению с заданием~2a добавлен класс \texttt{datasql} как ещё один наследник \texttt{data}.

% ============================================================
% ВСТАВЬТЕ СЮДА UML ЗАДАНИЯ 3:
% 1) Сгенерируйте на plantuml.com (скрипт diagram3.puml)
% 2) Сохраните как report/diagram3.png
% ============================================================
\begin{figure}[ht]
\centering
\includegraphics[width=0.95\textwidth]{diagram3.png}
\caption{Диаграмма классов для задания 3}
\label{fig:uml3}
\end{figure}

Скрипты PlantUML: \texttt{diagram3.puml} (классы), \texttt{diagram3\_er.puml} (ER).

\subsection{Выводы}

В задании~3 реализован класс \texttt{datasql} для чтения и записи каталога космических миссий в SQLite. Проверена цепочка обмена XML~$\to$~SQLite~$\to$~XML с использованием классов предыдущих заданий. Выполнены типовые запросы, построены диаграммы классов и ER.

% =====================================================================
% ЗАДАНИЕ 4
% =====================================================================
\clearpage
\section{Задание 4. Графический интерфейс на PyQt5}

\subsection{Постановка задачи}

С использованием библиотеки для программирования графического интерфейса Qt сделать форму для вывода данных в табличной форме, хранящихся в форматах XML, JSON или в базе данных SQLite, с использованием функций из заданий~2 и~3.

\subsection{Структура интерфейса}

Главное окно содержит вкладки:
\begin{itemize}[leftmargin=*]
\item \textbf{миссии} --- таблица миссий и форма редактирования (экипаж, агентство, год, длительность, эмблема);
\item \textbf{космонавты} --- таблица и форма ФИО;
\item \textbf{агентства} --- таблица и форма названия / краткого названия.
\end{itemize}

На каждой вкладке: таблица сверху, форма редактирования снизу, кнопки <<Добавить>>, <<Изменить>>, <<Удалить>>. Меню \texttt{File} позволяет создать новый каталог, открыть и сохранить данные в XML, JSON и SQLite.

\subsection{Базовые классы GUI}

Класс \texttt{rowCode} хранит соответствие номеров строк виджета и кодов записей. Класс \texttt{libWidget} содержит ссылку на каталог \texttt{cosmos}. Класс \texttt{dbTableWidget} наследуется от \texttt{QTableWidget} и \texttt{libWidget}, заполняет таблицу через переопределяемый метод \texttt{setData()}.

\begin{lstlisting}[caption={Фрагмент \texttt{dbTableWidget}}]
class dbTableWidget(QTableWidget, libWidget):
    def update(self, code=0):
        self.clearContents()
        self.setData()
        self.resizeColumnsToContents()
        self.resizeRowsToContents()
        if code:
            self.setCurrentCode(code)

    def setData(self):
        pass
\end{lstlisting}

Таблицы сущностей: \texttt{missionsTable}, \texttt{astronautsTable}, \texttt{agenciesTable}.

\subsection{Формы редактирования}

Базовый класс \texttt{editForm} задаёт сетку полей. Специализации:
\texttt{missionEditForm}, \texttt{astronautEditForm}, \texttt{agencyEditForm}.
Для миссии используются \texttt{astronautListWidget} (экипаж), \texttt{astronautCombo} (добавление в экипаж) и \texttt{agencyCombo} (выбор агентства).

Страница \texttt{page} объединяет таблицу, форму и кнопки; классы \texttt{missionPage}, \texttt{astronautPage}, \texttt{agencyPage} собирают конкретные виджеты. Класс \texttt{tabWidget} размещает страницы на вкладках.

\subsection{Главное окно}

\begin{lstlisting}[caption={Запуск приложения \texttt{main4.py}}]
from PyQt5.QtWidgets import QApplication
from mainwindow import mainWindow
import sys

app = QApplication(sys.argv)
mw = mainWindow()
mw.show()
sys.exit(app.exec_())
\end{lstlisting}

Класс \texttt{mainWindow} создаёт объект \texttt{cosmos} и объекты \texttt{dataxml}, \texttt{datajson}, \texttt{datasql}. Действия меню вызывают \texttt{readFile}/\texttt{writeFile} и обновляют вкладки через \texttt{tabWidget.update()}.

Запуск (при установленном PyQt5):
\begin{lstlisting}[language={},numbers=none]
python3 main4.py
\end{lstlisting}
Далее: \texttt{File $\to$ Open XML} и выбор файла \texttt{old.xml} (или JSON/SQLite).

\subsection{Диаграмма классов UML}

Диаграмма классов GUI приведена на рисунке~\ref{fig:uml4}.

% ============================================================
% ВСТАВЬТЕ СЮДА UML ЗАДАНИЯ 4:
% diagram4.puml -> report/diagram4.png
% ============================================================
\begin{figure}[ht]
\centering
\includegraphics[width=0.95\textwidth]{diagram4.png}
\caption{Диаграмма классов для задания 4}
\label{fig:uml4}
\end{figure}

\subsection{Выводы}

В задании~4 реализован графический интерфейс на PyQt5 для просмотра и редактирования каталога космических миссий с загрузкой и сохранением данных в XML, JSON и SQLite через классы предыдущих заданий.

\end{document}
